\documentclass[10pt,a4paper]{article}
\usepackage[T1]{fontenc}\usepackage[utf8]{inputenc}\usepackage[english,italian]{babel}
\usepackage{lmodern,microtype}\usepackage[a4paper,margin=1.85cm,headheight=15pt]{geometry}
\usepackage{enumitem,tabularx,array,longtable,booktabs,xcolor,hyperref,xurl,fancyhdr,titlesec,framed}
\definecolor{ddtadark}{RGB}{28,38,52}\definecolor{ddtablue}{RGB}{42,88,135}\definecolor{ddtagray}{RGB}{244,246,248}\definecolor{ddtagreen}{RGB}{48,111,78}\definecolor{ddtaorange}{RGB}{148,91,25}
\hypersetup{colorlinks=true,linkcolor=ddtablue,urlcolor=ddtablue}\pagestyle{fancy}\fancyhf{}\lhead{DDTA - BA0-R}\rhead{REVIEW-ONLY}\cfoot{\thepage}
\titleformat{\section}{\large\bfseries\color{ddtadark}}{}{0pt}{}\setlist[itemize]{leftmargin=1.45em,itemsep=.18em,topsep=.2em}
\setlength{\parindent}{0pt}\setlength{\parskip}{.32em}\setlength{\emergencystretch}{2em}
\newenvironment{factbox}[1]{\def\FrameCommand{\fcolorbox{ddtablue}{ddtagray}}\MakeFramed{\advance\hsize-2\width\FrameRestore}\noindent\textbf{#1}\par\smallskip}{\endMakeFramed}
\newenvironment{challenge}[1]{\def\FrameCommand{\fcolorbox{ddtaorange}{ddtagray}}\MakeFramed{\advance\hsize-2\width\FrameRestore}\noindent\textbf{#1}\par\smallskip}{\endMakeFramed}

\begin{document}
\begin{center}
{\LARGE\bfseries DDTA - Scheda di lettura compilata}\par
{\large SRC-0042 - BA0-R systems-modeling prior art}\par
{\small REVIEW-ONLY: da approvare prima del caricamento nel repository.}\par
\end{center}
\begin{factbox}{Identita bibliografica e accesso}
\begin{tabularx}{\linewidth}{>{\bfseries}p{3.4cm}X}
Source ID & SRC-0042\\
Titolo & NASA Systems Engineering Handbook, Revision 2\\
Autori/ente & Steven R. Hirshorn; Linda D. Voss; Linda K. Bromley; National Aeronautics and Space Administration (NASA)\\
Anno & 2017\\
Identificatore & NASA/SP-2016-6105 Rev 2; NTRS 20170001761\\
Accesso & internet\_public\_official\_nasa\_pdf\\
Verification state & verified\_full\_text\_public\\
\end{tabularx}
\end{factbox}
\section{1. Research problem and contribution}
Provides NASA lifecycle systems-engineering guidance across design, realization and technical management. It treats architecture as an interrelated structure of parts, relationships and connections and identifies multiple representational entities rather than one universal graph edge.
\section{2. Starting artifacts and generated representations}
Logical decomposition links functions and relationships to requirements before candidate solution selection. Architecture may be represented by functions, functional flows, interfaces, relationships, resource-flow items, physical elements, containers, modes, links and communication resources. Operational scenarios include events, actions, stimuli, information and interactions.
\section{3. Traceability, change, automation and human role}
Interface management explicitly identifies internal/external interfaces, origin, destination, stimuli and characteristics, then captures interface decisions, rationale, assumptions, anomalies and approved changes. Interfaces evolve as architecture emerges and are configuration-managed. Technical data management controls work products, decisions, assumptions, versions and access.
\section{4. Findings, limitations and relationship with DDTA}
The “Boundary of Technical Effort” is strong non-security prior art for a responsibility/control boundary: inside is controlled by the technical team, outside influences/is influenced but is not controlled. This supports distinguishing responsibility/system boundary from trust boundary. Architecture terminology is intentionally broad and practice-oriented, so it does not by itself establish a minimal ontology.
\section{Evidenza utile per BA0/BA1}
\begin{longtable}{p{1.2cm}p{2.0cm}p{4.0cm}p{8.0cm}}
\toprule
Pag. & Ruolo & Sezione & Interpretazione DDTA\\
\midrule
\endhead
62 & foundation & 4.3 Logical Decomposition & BA can contain implementation-independent system semantics grounded in requirements.\\
\addlinespace[.3em]
104 & foundation & 5.3 Product Verification Guidance & Behavioral representation needs more than static structure.\\
\addlinespace[.3em]
105 & foundation & Figure 5.3-2 End-to-End Data Flow & Flow is a distinct analyzable relation across heterogeneous structural elements.\\
\addlinespace[.3em]
145 & foundation & 6.3 Interface Management & Responsibility and interface boundaries are generic engineering concerns.\\
\addlinespace[.3em]
146 & foundation & 6.3.1.2 Interface Management Activities & Interface identity/change deserves explicit consideration rather than being inferred only from interaction text.\\
\addlinespace[.3em]
147 & foundation & 6.3.1.3 Outputs & Architecture/interface provenance and change history are established engineering needs.\\
\addlinespace[.3em]
186 & challenge & Appendix B Architecture & Actor+Asset+Interaction is far too small as a generic system-modeling vocabulary.\\
\addlinespace[.3em]
236 & foundation & Appendix J 3.5 Boundary of Technical Effort & Strong evidence for responsibility/control boundary distinct from security trust.\\
\addlinespace[.3em]
\bottomrule
\end{longtable}
\section{Bibliography mining / follow-up}
\begin{itemize}
\item NASA-HDBK-1009A
\item ISO/IEC/IEEE 15288
\item INCOSE Systems Engineering Handbook
\end{itemize}
\begin{challenge}{Governance note}
Questa scheda e un ausilio di review. Le future citazioni di tesi devono risolversi alla fonte registrata e a una posizione esatta nell'excerpt ledger approvato. Nessun concetto esterno diventa automaticamente un BAE DDTA.
\end{challenge}
\end{document}
